\chapter{Answer Set Programming}

\section{Introduzione}

Durante la cosiddetta "War of Semantics" nasce l'esigenza di dare una semantica alla negazione per fallimento adottata dagli interpreti PROLOG.

\dfn{Answer Set Programming}{
  Paradigma di programmazione in cui le soluzioni sono i \newfancyglitter{modelli} (Answer Set), non più le prove
}

\nt{Con PROLOG ha in comune solo la sintassi, per il resto è tutt'altra cosa.}

\paragraph{L'Answer Set Programming (ASP):}

\begin{itemize}
  \item È particolarmente utile per risolvere problemi combinatori (soddisfacimento di vincoli, planning).
  \item ASP solvers molto efficienti sviluppati per supportare questa metodologia (DLV, smodels, \fancyglitter{CLINGO}, Cmodels,\dots).
\end{itemize}

\nt{Un ASP solvers è l'equivalente di un interprete PROLOG.}

\paragraph{Codice ASP:}

\begin{itemize}
  \item Insieme finito di regole: $a :- b_1, b_2, \dots, b_n, not c_1, not c_2, \dots, not c_m$. 
  \item $a, b_i, c_j$ sono letterali nella forma $p$ on $- p$:
    \begin{itemize}
      \item - è la negazione classica. 
      \item $not$ è la negazione per fallimento.
    \end{itemize}
  \item $a$ è opzionale, senza si ha \fancyglitter{integrity constrain} (regole senza testa): 
    \begin{itemize}
      \item $:- a_1, a_2, \dots, a_k$.
      \item È inconsistente che siano tutti veri\dots
      \item Serve per filtrare/buttare via dei modelli.
    \end{itemize}
  \item Si applica ai soli programmi logici proposizionali. 
  \item La maggior parte dei tool per ASP consente per
comodità di usare variabili, ma le clausole devono poter
essere trasformate in un numero finito di clausole ground.
\end{itemize}

\paragraph{ASP vs. PROLOG:}

\begin{itemize}
  \item In ASP l’ordine dei letterali non ha alcuna importanza. 
  \item Prolog è goal-directed, ASP no. 
  \item La SLD-risoluzione del Prolog può portare a loop,
mentre gli ASP solver non lo consentono. 
\item PROLOG ha il cut(!), ASP no.
\end{itemize}

\subsection{Negazione}

\begin{itemize}
  \item \fancyglitter{Classica}: 
    \begin{itemize}
      \item attraversa :- treno. 
      \item Si attraversa solo se si può derivare che il treno non è in
arrivo.
    \end{itemize}
  \item \fancyglitter{Per fallimento}: 
    \begin{itemize}
      \item attraversa :- not treno. 
      \item Si può attraversare in assenza di informazione esplicita sul treno in
arrivo.
    \end{itemize}
  \item Un letterale negato -p non ha nessuna proprietà particolare. 
  \item Viene considerato come se fosse un nuovo atomo positivo,
aggiungendo il vincolo :- p, -p.
\end{itemize}

\clm{}{}{
  CLINGO: 
  \begin{itemize}
    \item Fornisce modelli e indica se sono \fancyglitter{SATISFIABLE} (soddisfacibili) o \fancyglitter{UNSATISFIABLE} (insoddisfacibili). 
    \item Se si aggiunge il parametro "0" vengono mostrati tutti i modelli (ATTENZIONE: potrebbero essere migliaia, è sconsigliato metterlo di default).
  \end{itemize}
}

\section{Semantica}

\dfn{Answer Set}{
  Un Answer Set è un modello minimale (stabile).
}

\nt{Un programma ASP privo di letterali $not p_i$ ha un unico modello minimale che è il suo answer set.

Potrebberò esserci più answer set, ma interessa solo quello minimale.}

\qs{}{Che succede se è presente la negazione per fallimento?}

\dfn{Ridotto}{
  Il \newfancyglitter{ridotto} $P^S$ rispetto a un insieme di atomi S: 
  \begin{itemize}
    \item Rimuove ogni regola il cui corpo contiene $not L$, per $L \in S$. 
    \item Rimuove tutti i $not L$ dai corpi delle restanti regole.
  \end{itemize}

  $P^S$ non contiene atomi con negazione per fallimento: 
  \begin{itemize}
    \item Ha un unico answer set. 
    \item Se tale answer ser coincide con S, allora S è un answer set per P.
  \end{itemize}
}

\nt{In ASP si può usare \#show/cardinalità per mostrare solo alcune parti dei modelli.}
\begin{figure}[h]
    \centering
    \includegraphics[scale=0.2]{03/show.png}
    \caption{"Show" or something idk.}
\end{figure}

